\documentclass[../Project-Report-Main.tex]{subfiles}

\begin{document}

\chapter{Literature Review}

\section{Overview of Personal Finance Management (PFM) Systems}

Personal Finance Management (PFM) systems have evolved as fundamental digital tools aiding individuals in tracking cash flows, formulating budgets, managing debt, and planning savings. Early literature in consumer economics emphasizes that active monitoring of personal financial transactions significantly improves financial awareness, reduces impulse spending, and improves personal savings rates.

Historically, PFM tools were classified into three main generations:
\begin{enumerate}
 \item \textbf{First Generation (Manual Bookkeeping):} Paper ledgers, physical envelope budgeting systems, and offline spreadsheet tools (Microsoft Excel).
 \item \textbf{Second Generation (Desktop \& Early Web Software):} Intuit Quicken, Microsoft Money, and downloadable desktop applications requiring manual input or file imports (QIF/OFX format).
 \item \textbf{Third Generation (Cloud-Based \& Mobile PFM Apps):} Mint, YNAB (You Need A Budget), PocketGuard, and mobile apps offering automatic bank syncing or SMS tracking.
\end{enumerate}

Recent academic studies highlight that despite the availability of third-generation apps, adoption barriers persist due to data privacy concerns, complex setup processes, and inability to handle region-specific financial formats (such as Indian UPI payment logs and password-protected PDF statements).

\section{Evolution of Transaction Tracking: Manual vs. Automated Methods}

Literature comparing manual and automated expense tracking methodologies highlights a stark trade-off between user friction and data accuracy:

\begin{itemize}
 \item \textbf{Manual Tracking Systems:}
 \begin{itemize}
   \item \textit{Strengths:} Total user privacy, high awareness during entry, no integration required.
   \item \textit{Weaknesses:} High cognitive load, prone to human error, omission of micro-transactions, and high abandonment rates (over 70\% within 30 days according to PFM behavioral studies).
 \end{itemize}

 \item \textbf{Automated Tracking via SMS Scraping:}
 \begin{itemize}
   \item \textit{Strengths:} Low user friction, real-time transaction detection.
   \item \textit{Weaknesses:} Intrusive permissions (requires full access to personal SMS inbox), inability to capture details when SMS notification fails, lack of support for detailed itemized statements, and restriction to Android mobile platforms (iOS restricts SMS reading).
 \end{itemize}

 \item \textbf{Automated Tracking via PDF Statement Parsing:}
 \begin{itemize}
   \item \textit{Strengths:} High fidelity, captures official bank records, supports multi-line transaction narratives, works cross-platform (web and mobile), respects user privacy by operating on user-uploaded files without background message scanning.
   \item \textit{Weaknesses:} Format variability across financial institutions, password encryption, requirement of specialized PDF parsing logic.
 \end{itemize}
\end{itemize}

The literature concludes that PDF statement parsing provides the ideal middle ground for users seeking privacy-conscious, high-fidelity automated tracking.

\section{Bank Statement Ingestion \& Parsing Techniques}

Document processing literature divides electronic bank statement parsing into three primary techniques:

\begin{enumerate}
 \item \textbf{Optical Character Recognition (OCR):} Tools like Tesseract OCR or AWS Textract scan statement images/PDFs. While effective for scanned paper documents, OCR requires heavy processing power, is prone to character recognition errors (e.g., mistaking '0' for 'O' or '1' for 'l'), and is unnecessarily computationally expensive for native digital PDFs.

 \item \textbf{Table Extraction Frameworks (e.g., Tabula, Apache PDFBox):} Frameworks designed to extract raw text and tabular structures directly from native digital PDF streams. Apache PDFBox is widely acknowledged in open-source software literature for its low memory footprint, speed, and robust handling of PDF document security, including decryption via user password parameters.

 \item \textbf{Regex-Driven Tokenization:} In financial domain applications, raw extracted text streams must be converted into structured database records. Research demonstrates that combining \texttt{PDFTextStripper} output with deterministic Regular Expressions (Regex) allows efficient extraction of date markers, transaction types, debit/credit fields, and multi-line narrations even when tabular gridlines are missing from the PDF structure.
\end{enumerate}

\section{Categorization Engines: Rule-Based vs. Machine Learning Approaches}

Categorizing unorganized financial transactions is a central research topic in financial informatics. Existing literature analyzes two dominant paradigms:

\begin{enumerate}
 \item \textbf{Machine Learning (ML) Classifiers:} Supervised learning algorithms (such as Naive Bayes, Support Vector Machines, and Random Forests) and Natural Language Processing (NLP) models (TF-IDF, word embeddings) have been employed to classify transaction strings. While effective at scale, ML models require large, pre-labeled training datasets, suffer from cold-start problems for new users, and can misclassify unique merchant names without continuous retraining.

 \item \textbf{Rule-Based \& Keyword Mapping Engines:} Rule-based systems apply deterministic pattern matching (e.g., checking if narration contains "SWIGGY" or "PETROL"). Literature indicates that rule-based engines achieve near 100\% precision for known merchants and allow instantaneous user customization.

 \item \textbf{Hybrid Categorization Architecture:} Recent literature strongly advocates a hybrid approach: leveraging persistent user-defined rules as high-priority overrides, combined with static keyword dictionaries for general merchants, and fallback default categories. This hybrid architecture eliminates training overhead while maintaining high accuracy and user customization—an approach directly adopted in Finlytics.
\end{enumerate}

\section{Integration of Generative Artificial Intelligence (AI) in Financial Advisory}

The emergence of Large Language Models (LLMs) such as OpenAI GPT-4, Meta LLaMA 3, and Anthropic Claude represents a paradigm shift in financial technology. Traditional PFM systems relied on hardcoded heuristic logic to output static tips. In contrast, LLM-driven advisory engines offer natural language synthesis, contextual understanding, and personalized financial coaching.

Key findings from recent AI in Fintech literature include:
\begin{itemize}
 \item \textbf{Context Aggregation:} LLMs perform exceptionally well when provided with aggregated metadata summaries (total income, expenses, category spending, top merchants) rather than raw individual transaction logs, preserving privacy while enabling deep insights.
 \item \textbf{LLaMA 3.1 Architecture:} Meta's LLaMA 3.1 8B Instant model (accessed via high-speed inference platforms like Groq API) provides fast inference times ($< 1$ second) and high reasoning capabilities suitable for real-time web applications.
 \item \textbf{Latency \& Cost Optimization:} Literature emphasizes the necessity of prompt sanitization and session caching strategies (e.g., caching insight responses for 2--5 minutes) to avoid redundant API calls, control computational overhead, and deliver responsive user experiences.
\end{itemize}

\section{Comparative Study / Literature Comparison Matrix}

Table~\ref{tab:lit_comparison} illustrates a detailed comparison of existing personal finance management approaches against the proposed Finlytics system across major functional and technical parameters.

\begin{table}[htbp]
  \centering
  \caption{Literature Comparison Matrix of PFM Approaches}
  \label{tab:lit_comparison}
  \small
  \begin{tabular}{|p{1.4in}|p{1.0in}|p{1.1in}|p{1.1in}|p{1.3in}|}
    \hline
    \textbf{Feature / Dimension} & \textbf{Manual Tracker (Excel/Sheets)} & \textbf{SMS-Based Apps (e.g., Walnut)} & \textbf{Generic PDF Converters} & \textbf{Proposed System (Finlytics)} \\
    \hline
    Automated PDF Parsing & No & No & Partial (Text) & Yes (Full) \\
    \hline
    Password PDF Support & No & No & No & Yes (SBI YONO) \\
    \hline
    Privacy / Data Security & High & Low (Scans SMS) & Medium & High (Isolated DB) \\
    \hline
    Smart Categorization & Manual & Automated & None & Hybrid Engine \\
    \hline
    Persistent User Rules & Manual & Limited & None & Yes (MerchantRule) \\
    \hline
    Budget Overspend Alerts & Manual formulas & Yes & No & Yes (Visual progress) \\
    \hline
    LLM Financial Advice & No & No & No & Yes (Groq/LLaMA 3.1) \\
    \hline
    Cross-Platform Web UI & Yes & Mobile Only & Web Only & Yes (Full-Stack) \\
    \hline
    Open Source / Self-Host & Yes & Proprietary & Variable & Yes \\
    \hline
  \end{tabular}
\end{table}

\section{Identified Research Gaps \& Rationale for Finlytics}

Based on the literature review and comparative analysis, the following research and technological gaps were identified:

\begin{enumerate}
 \item \textbf{Lack of Integrated Support for Indian Bank PDF Formats:} Most global or open-source PFM tools lack dedicated parsing patterns for Indian banking statements (such as password-protected SBI YONO PDFs and multi-line Google Pay UPI narratives).
 \item \textbf{Privacy Vulnerabilities in Automated Mobile Trackers:} Popular automated mobile tracking applications compromise user privacy by requiring full SMS access permissions.
 \item \textbf{Absence of Actionable AI Advisory in Lightweight Web PFM Systems:} Existing lightweight or self-hosted PFM web tools provide charts but lack generative AI capabilities to analyze 30-day financial behaviors and output natural language recommendations.
\end{enumerate}

\noindent \textbf{Rationale for Finlytics:} \\
Finlytics directly addresses these research gaps by synthesizing password-protected PDF parsing (Apache PDFBox), a hybrid merchant categorization engine, interactive visual analytics (Chart.js), and generative LLM financial advisory (Groq LLaMA 3.1) into a single, secure, full-stack Spring Boot web application.

\end{document}
